%%% ANNOTATED MARKUP DRAFT — read this block, then delete it before compiling for real.
%
% Base file: e_methodology_chapter_draft_14th_aug_6_29.tex (confirmed by Shreeya as the current
% ~18-page working draft, 2026-08-14).
%
% What this file is: a TRACKED-CHANGES version, not a clean final chapter. Original prose before
% the "RQ3: Plot-specific growth-curve deviation" section is preserved verbatim; every proposed
% cut, move, or restructure is flagged inline with \ai{...} (brown text) rather than silently
% applied, per the editing brief. Where a proposed table is a genuine space-saving replacement
% for prose, the real replacement table is given immediately after its \ai{} note (not colored —
% it's usable content, not an instruction) so you can drop it in directly once you accept the cut.
% From the RQ3 section onward, the brief allows more direct rewriting (that material "has not yet
% been closely edited"), so it is rewritten cleanly with plain HTML-style % comments marking the
% larger structural moves (e.g. results-statements relocated to the Results chapter) rather than
% \ai{} wrapping every sentence.
%
% Companion deliverable: the chat response accompanying this file contains the full section-by-
% section editing plan (A-G): diagnosis, per-section classification + page-saving estimates,
% revised ~12-page structure, figure/table specs, equation audit, limitations rewrite, and the
% results-leakage audit. Read that alongside this file — this file is the "what", that is the
% "why" and the running total of pages saved.
%
% Correctness bugs fixed directly rather than only \ai-flagged (structural/LaTeX bugs, not
% stylistic edits, so they don't need the "preserve wording, mark every change" treatment — noted
% here as a single list instead of individually):
% (1) tab:feature-sets-composition-rq1 had two \caption commands in one float (source lines 142
%     and 147); the first, less-polished one is removed.
% (2) tab:feature-set-purpose's tabular declared 5 columns but every row supplied 6 fields (a
%     "comparison role" column, e.g. "Primary"/"Secondary", had no header cell); a 6th column is
%     added to both the header and the column spec.
% (3) RQ2b's heading was an unlabelled \section* (unnumbered, dropped from the ToC) while every
%     sibling section is a numbered, labelled \section — changed to match, with \label{sec:rq2b}
%     added, since other cross-references in the chapter already assumed it existed.
% (4) Several \S\ref{...} cross-references in the source point at labels that were never defined
%     anywhere in the chapter (sec:feature-final, sec:rq1-pinn-loss, sec:rq1-ablation) — these
%     would render as "??" in a compiled PDF. Labels are added at the section/subsection/paragraph
%     they evidently meant. sec:common-design and sec:rq1-models are also newly labelled, since
%     this rewrite's own added cross-references point at them.
% All four are pointed out again at their location below where relevant.
%%%

%%% Colour commands - delete for final draft
% --- Editing Macros (Hierarchy of Importance) ---
% 1. FACT CHECK (Critical): Red commands immediate attention for potential inaccuracies.
\newcommand{\fc}[1]{\textcolor{red}{[FactC: #1]}}
% 2. CITE (High): Orange highlights missing evidence or attribution.
\newcommand{\ct}[1]{\textcolor{orange}{[CITE: #1]}}
% 3. APPENDIX (Medium): Blue indicates structural moves (shifting text out of the main body).
\newcommand{\app}[1]{\textcolor{blue}{[APP: #1]}}
% 4. SMALL REWORD (Low): Teal is for minor stylistic edits, mirroring your friend's setup.
\newcommand{\sr}[1]{\textcolor{teal}{[Sreword: #1]}}
% 5. FIGURE (Low): Purple is for figure-related notes (refer to figure, create figure, etc.).
\newcommand{\fig}[1]{\textcolor{purple}{[FIG: #1]}}
% 6. CLAUDE EDIT (this pass): Brown, per the editing brief. Used for every proposed change before
% the RQ3 section (deletions, replacements, moves, table/figure conversions), and for factual
% corrections anywhere in the chapter.
\newcommand{\ai}[1]{\textcolor{brown}{[AI: #1]}}

%=======================================================================================
\chapter{Methodology and Experimental Design}
\label{ch:methodology}

This chapter outlines how the raw data was utilised to answer the three research questions. It is organised by method, first the the experimental framework, then common design decisions and data leakage control, the external environmental feature set construction, the top-height prediction models, the two attribution analyses, evaluation design, and finally assumptions, limitations and reproducibility.
\ai{Fix typos only: "first the the experimental" -> "first the experimental". No content change; the roadmap sentence is accurate and Table~\ref{tab:rq-framework} immediately below already restates it visually, so shortening further would just create the same duplication the rest of this pass is trying to remove.}

%------------------------------------------------
\section{Research questions and experimental framework}
\textbf{RQ1} test which models accurately predict top height and under what input and cohort year. \textbf{RQ2} is an attribution task, asking which areas over or under predict top height relative to the the average CR equation formed from the entire dataset. \textbf{RQ3} fits each plot with its own asymptote from only the plots repeated observations, and makes comparisons on the growth trajectories of curves, and if they are associated with specific stand, management or external environmental variables.
\ai{Grammar/typo pass only, no content change: "RQ1 test" -> "RQ1 tests"; "relative to the the average" -> "relative to the average"; "the plots repeated observations" -> "the plot's own repeated observations". Keep the rest verbatim -- it is already the shortest accurate statement of what each RQ does, and Table~\ref{tab:rq-framework} carries the detail.}

\begin{table}[!htbp]
\scriptsize
\renewcommand{\arraystretch}{1.2}
\begin{tabular}{p{0.8cm}p{1.4cm}p{3.3cm}p{4.2cm}p{2.9cm}}
\toprule
RQ & Unit & Target & Model families compared & Evaluation design \\
\midrule
RQ1 & Plot-survey row & 95th-percentile LiDAR top height & Chapman-Richards curve,
linear/random forest/XGBoost, deep neural network, Chapman-Richards-guided neural network
variants & Spatial-block (primary), pooled spatial $k$-fold, temporal, plot-level \\
RQ2 & Plot & Mean residual from the single pooled Chapman--Richards curve, across a plot's survey
history & Elastic Net, XGBoost, two-stage mixed-effects model & Spatial-block \\
RQ3 & Plot & Difference between a plot's own fitted asymptote and its yield-class-implied
asymptote & Elastic Net, XGBoost, GNNWR & Spatial-block \\
\bottomrule
\end{tabular}
\caption{Research questions mapped to unit of analysis, target, model families, and evaluation
design.}
\label{tab:rq-framework}
\end{table}
\ai{Keep this table unchanged -- it is already the most space-efficient statement of the framework and everything after it in the chapter should refer back to it rather than restate its columns in prose.}

\ai{RELOCATED (2026-08-14, after discussion): this pipeline-overview figure originally sat at the end of \S\ref{sec:feature-final}, right before RQ1 -- moved here instead, as the chapter's second figure (after Table~\ref{tab:rq-framework}), on your call to keep and reposition it rather than cut it. Kept deliberately, against the general space-reduction instinct, because a flowchart of cohorts $\to$ splits $\to$ feature screening $\to$ three branches $\to$ outputs does a job Table~\ref{tab:rq-framework} cannot: it shows the reader the pipeline's \emph{shape} before any box is individually unpacked, which is exactly the value a diagram has over prose in an ML methodology chapter. Its box labels (e.g.\ "Set1-4 feature screening", "spatial and temporal partitions") don't need to be fully defined yet at this point in the chapter -- that's normal for an opening overview figure, and every term is properly defined in the sections that follow. Per your instruction, all of the mechanical detail that used to live in a five-part comment block (a prose caption draft, a "confirmed outputs" note, a mermaid sketch, and a prose fallback -- four overlapping descriptions of the same thing) is now consolidated into the one real caption below, and nothing describing the pipeline is left duplicated in body prose elsewhere in the chapter.}

\ai{PAGE-BUDGET NOTE: keeping this figure removes the $\sim$0.4-0.5 page saving it was counted for in the original 12-page plan (see part C of the review). Restoring it alone puts the chapter closer to $\sim$13-13.5 pages rather than 12. Two places I'd look first if you want to claw that back, rather than me cutting them unilaterally: (1) Table~\ref{tab:feature-set-purpose}'s "Purpose" column is still full sentences and could tighten to phrases; (2) this section's own RQ1 opening paragraph (\S\ref{sec:rq1-models}) carries a parenthetical aside about stakeholder investment decisions that isn't load-bearing for the methodology itself. Flagging both, not cutting either -- your call.}

% FIGURE PLACEHOLDER -- pipeline overview. Content: a flowchart, cohorts feeding data-splitting
% (Table~\ref{tab:splits-summary}) and, separately, the shared Set1-4 screening process
% (Table~\ref{tab:feature-sets-composition-rq1}); both outputs feed three parallel branches, one
% per research question (RQ1 top-height prediction, RQ2 shared-curve residual attribution, RQ3
% plot-specific asymptote attribution); each branch terminates in its own confirmed evaluation
% output (see caption). Do not draw a rendered coefficient/residual map as an output box -- not
% confirmed to exist in the modelling pipeline as of this pass.
\begin{figure}[!htbp]
\centering
\fbox{\parbox{0.9\textwidth}{\centering\vspace{4cm}FIGURE PLACEHOLDER --- pipeline overview, see
comment above source for exact content\vspace{4cm}}}
\caption{Overview of the full pipeline from raw data to each research question's results. Both
cohorts (\S\ref{sec:data-cohorts}) feed the four data-splitting strategies
(Table~\ref{tab:splits-summary}) and a shared screening process that produces each research
question's own Set1--4 (Table~\ref{tab:feature-sets-composition-rq1}). Three branches follow, one
per research question (Table~\ref{tab:rq-framework}): RQ1's top-height prediction models, RQ2's
shared-curve residual attribution, and RQ3's plot-specific asymptote attribution. RQ1 produces
per-fold and pooled predictive metrics under each split type; RQ2 produces a fixed-effects
coefficient table plus held-out residual predictions; RQ3 produces per-fold predictive metrics and
a per-plot coefficient table (GNNWR), together with a residual spatial-autocorrelation diagnostic
(\S\ref{sec:eval-diagnostics}). No branch currently produces a rendered coefficient or residual
map.}
\label{fig:pipeline-overview}
\end{figure}

%------------------------------------------------
\section{Common Experimental Design and Data Leakage control}
\label{sec:common-design}

\paragraph*{Evaluation Framework and Data Splits}. 4survey is the main cohort for every research question, used for the spatial-generalisation splits; 6survey (\S\ref{sec:data-cohorts}) is used for the temporal split and run in parallel elsewhere as a secondary check wherever sample size permits. For RQ1, each row in the dataset represents a single plot in a given survey year. For RQ2 and RQ3, data are aggregated so each row represents a single plot (Table~\ref{tab:rq-framework}). Model performance is evaluated using four distinct data-splitting strategies:
\ai{This paragraph's first sentence re-explains the 4survey/6survey cohort roles, which Chapter~\ref{ch:data} \S\ref{sec:data-cohorts} already states in full (the nested-cohort relationship, why 4survey is primary, why 6survey serves the temporal question). DELETE the first sentence and REPLACE with: "As established in \S\ref{sec:data-cohorts}, 4survey is the primary cohort throughout; 6survey's longer per-plot trajectories suit the temporal split, and its spatial-block results are compared against 4survey's throughout to help explain differences in model behaviour." (revised 2026-08-15: the earlier draft of this replacement framed 6survey's spatial-block role as only a "secondary check, where sample size allows," which undersold it -- 6survey is genuinely run for spatial-block comparisons throughout, and that cross-cohort comparison is itself an analytical tool, e.g.\ explaining why DNN inflates more than PINN under easy splits or why GNNWR's stability differs by cohort, not just a backup confirmation.) Keep the rest of the paragraph (RQ1 vs RQ2/RQ3 row unit, and the lead-in to the four splits) as-is.}

\begin{itemize}
\setlength\itemsep{0pt}
\item \textbf{Plot-level}. Individual plots are split at random (60/20/20, train/validation/test, seed 42), following previous work \ct{source that uses plot level split}. However a held out plots nearest neighbour is typically a training plot meters away, so a model can succeed by picking on its location in Aberfoyle rather than learning from features that shape growth trajectory. Kept as a baseline reference for comparison.

\item \textbf{Spatial block} (primary). Addresses spatial data leakage by assigning entire forestry compartments to train/validation/test sets (60/20/20, row-count-balanced, seed 42). To eliminate edge effects, any training plot within 60\,m of a held-out plot is excluded (validation/test plots are always retained; Figure~\ref{fig:split-maps}, bottom row). Primary isolation is achieved at the compartment level, while the 60\,m buffer—based on typical plot spacing, reduces local spatial autocorrelation along boundaries without discarding excessive data. \citep{roberts2017Cross_C5_EVAL}.

\item {\textbf{Spatial block, pooled $k$-fold}} (stronger confirmation of the single spatial-block split). Compartments were assigned to five folds balanced approximately by row count. Each fold was held out once as the test set, the following fold was used for validation, and the remaining folds were used for training. Thus, every compartment received one prediction from a model that was not trained on that compartment. Held-out predictions from all five folds were pooled before calculating the overall performance metrics, avoiding distortions caused by differences in fold size and the range of observed heights\cite{ploton2020Spatial_C5_EVAL}.

\item \textbf{Temporal} (secondary, harder test). Whole survey years, not compartments, are assigned to train/validation/test: 4survey trains on 2008 and 2012, validates on 2021, tests on 2023; 6survey trains on 2002/2006/2008/2012, validates on 2021, tests on 2023. This asks whether a model trained on earlier growth still predicts a real future survey it never saw, a different, harder question than spatial generalisation (Figure~\ref{fig:temporal-split}), following on from previous work \cite{lynch2025Digital_C4_NN}.
\end{itemize}
\ai{This itemize is the single largest space-reduction opportunity before RQ3 -- four wrapped paragraphs to state what is fundamentally a small parameter table. DELETE this itemize (all four items) and REPLACE with two things: (1) one short paragraph carrying ONLY the two sentences that are genuine methodological argument, not parameters -- the plot-level leakage mechanism ("a held-out plot's nearest neighbour is typically a training plot metres away...") and the temporal split's different, harder question -- and (2) Table~\ref{tab:splits-summary} below for every construction detail (ratios, seed, buffer, fold count, years, citations). This keeps every fact, removes the wrapped-paragraph overhead of four bulleted blocks of prose.}

\ai{Replace the preceding itemize with the following paragraph + table (real content, not a placeholder -- drop in directly if you accept the cut):}

Nearby held-out plots at the level of individual plots leak location information -- a held-out plot's nearest neighbour is typically a training plot metres away, so a model can succeed by learning \emph{where} it is in Aberfoyle rather than \emph{why} a plot grows the way it does. Spatial block splitting removes this by holding out whole compartments instead. The temporal split asks a different, harder question: whether a model trained on past growth still predicts a real future survey it never saw, rather than an unseen location in the same time period. Table~\ref{tab:splits-summary} gives the exact construction of all four splits used across RQ1--RQ3 (Table~\ref{tab:rq-framework}).

\begin{table}[!htbp]
\centering
\scriptsize
\renewcommand{\arraystretch}{1.2}
\begin{tabular}{p{2.1cm}p{5.1cm}p{5.3cm}}
\toprule
Split & Construction & Role \\
\midrule
Plot-level & Individual plots, 60/20/20, random, seed 42 \ct{source that uses plot-level split}. & Baseline reference only, not a leakage-controlled split. \\
Spatial block (primary) & Whole compartments, 60/20/20, row-count-balanced, seed 42; training plots within 60\,m of a held-out plot excluded (val/test always retained) \citep{roberts2017Cross_C5_EVAL}. & Main leakage control for every RQ (Table~\ref{tab:rq-framework}). \\
Spatial block, pooled $k$-fold & 5 compartment folds, row-count-balanced; each held out once as test, next fold as validation, remainder as training; held-out predictions pooled before scoring \citep{ploton2020Spatial_C5_EVAL}. & Stronger confirmation of the single spatial-block split -- every compartment scored exactly once. \\
Temporal (secondary) & Whole survey years assigned to train/val/test: 4survey trains 2008+2012, validates 2021, tests 2023; 6survey trains 2002--2012, validates 2021, tests 2023 \citep{lynch2025Digital_C4_NN}. & Tests generalisation across time (RQ1 only; RQ2/RQ3 have no temporal axis, \S\ref{sec:eval-diagnostics}). \\
\bottomrule
\end{tabular}
\caption{The four data-splitting strategies. All use seed 42; compartment/sub-compartment
identifier columns are never included in any model's feature list, enforced automatically.}
\label{tab:splits-summary}
\end{table}

Every model uses the same split assignment for a given (cohort, split type, seed) combination, so results are always compared on identical train/validation/test membership. Compartment/block/sub-compartment identifier columns are never included in any model's feature list, enforced automatically. Shared data partitions and matched architectures reduce confounding when comparing model families under the same split e.g. the DNN and the Chapman-Richards-PINN. - though residual differences can still reflect optimisation behaviour or stochastic variation rather than the model design alone. %TODO: CHECK IF PINN ARCHTEICTURE IS SAME AS DNN.
\ai{The identifier-column sentence is now redundant with Table~\ref{tab:splits-summary}'s own caption (added above) -- DELETE it here. Also fix "ARCHTEICTURE" typo and resolve the open TODO before submission (verify PINN/DNN share an architecture; \S\ref{sec:rq1-pinn-loss} says they do, so this TODO is likely already answered by that section -- confirm and delete the comment, don't leave it in a submitted draft). Tighten the remaining sentence to: "Every model uses the same split assignment for a given (cohort, split type, seed) combination. Shared partitions and matched architectures (e.g. DNN vs.\ CR-PINN) reduce confounding when comparing model families under the same split, though residual differences can still reflect optimisation behaviour or stochastic variation rather than model design alone."}

\begin{figure}[!htbp]
\centering
\begin{minipage}[b]{0.32\textwidth}
  \centering
  \fbox{\parbox{\textwidth}{\centering\vspace{2.5cm}\textbf{Panel A}\\Plot-Level Split\vspace{2.5cm}}}
\end{minipage}\hfill
\begin{minipage}[b]{0.32\textwidth}
  \centering
  \fbox{\parbox{\textwidth}{\centering\vspace{2.5cm}\textbf{Panel B}\\Spatial-Block Split ($60\text{~m}$ Buffer)\vspace{2.5cm}}}
\end{minipage}\hfill
\begin{minipage}[b]{0.32\textwidth}
  \centering
  \fbox{\parbox{\textwidth}{\centering\vspace{2.5cm}\textbf{Panel C}\\Temporal Split Timeline\vspace{2.5cm}}}
\end{minipage}
\caption{Overview of data-splitting strategies across cohorts: (A) Unconstrained random plot-level split showing nearby spatial leakage; (B) Compartment-based spatial-block split with a $60\text{~m}$ boundary exclusion buffer; and (C) Sequential temporal split timeline across survey years.}
\label{fig:splits-combined}
\end{figure}
\ai{FIGURE: Show three panels left to right on real Aberfoyle geography (reuse whatever map base layer the data chapter's Figure~\ref{fig:data-cohort-map} uses, for visual consistency). Panel A (plot-level): a small cluster of plots coloured by train/val/test membership, with at least one held-out plot's nearest training-plot neighbour explicitly connected by a short line/arrow labelled with its distance in metres -- this is the one fact prose alone struggles to make visceral, so it must be visible, not just captioned. Panel B (spatial-block): compartment polygons coloured by train/val/test, with the 60\,m exclusion buffer shown as a hatched ring around held-out compartments, and a callout noting buffer-excluded training plots are dropped, not reassigned. Panel C: a horizontal timeline bar for each cohort (4survey: 2008, 2012 train -- 2021 val -- 2023 test; 6survey: 2002/2006/2008/2012 train -- 2021 val -- 2023 test), with the 2012-2021 gap visually marked. Caption must state: (1) entire compartments, not individual plots, are held out in B to control spatial leakage; (2) the buffer only removes training-side plots, val/test membership is never altered; (3) the exact train/val/test years per cohort shown in C. This figure replaces nothing further in the prose (the itemize above it is already cut back to two argument sentences) -- it is retained because a map is the only way to make the plot-level leakage mechanism and the buffer's asymmetry (train-only) legible at a glance, which Table~\ref{tab:splits-summary} cannot do.}

\paragraph*{Preprocessing and leakage prevention}. Continuous features are standardised to zero mean and unit variance using parameters fitted exclusively on the training partition to prevent scaling data leakage. Age is scaled independently using a separate scaler, allowing its training-split standard deviation to be extracted directly for the physics-informed loss function's unit conversion (\S\ref{sec:rq1-pinn-loss}??).( Unordered categorical soil variables are one-hot encoded to avoid imposing artificial numeric orders.  -- make this sentnece simpler an shosrter. )Any plot with missing values in required features is removed entirely across all survey years to ensure models never encounter incomplete temporal trajectories. For the terrain-conditioned RQ1 models this removes 39 plots (156 rows) from 4survey and none from 6survey. \app{add to appendix which cells these were or where they are located on the map}
\ai{Fix the broken "\S\ref{sec:rq1-pinn-loss}??" (stray double question mark -- looks like an unresolved cross-reference check, confirm the label exists and delete the "??"). Resolve your own inline note ("make this sentence simpler and shorter") and the stray unmatched parenthesis around the soil-encoding sentence. REPLACE the whole paragraph with: "Continuous features are standardised (zero mean, unit variance) using parameters fitted only on the training partition; age is scaled separately so its training-split standard deviation can be extracted for the physics-loss unit conversion (\S\ref{sec:rq1-pinn-loss}). Categorical soil variables are one-hot encoded. Any plot missing a required feature in any survey year is removed entirely, so no model ever sees an incomplete temporal trajectory -- 39 plots (156 rows) for the terrain-conditioned RQ1 models on 4survey, none on 6survey (affected plot IDs and map location: Appendix~\ref{app:missing-cells})." This resolves the \app{} note with an actual appendix pointer instead of leaving it as an open instruction to yourself.}

%------------------------------------------------
\section{Environmental feature-set construction}
\label{sec:feature-final}
Candidate environmental and management variables are screened separately for RQ1 (predictive tiers), RQ2 (shared-curve attribution), and RQ3 (plot-level attribution). All three follow the same recipe below, differing only in candidate pool and
target column. This produces four tiered feature sets (Set1–Set4; \S\ref{sec:feature-final}) to evaluate whether expanding environmental inputs improves model performance or introduces unhelpful noise across different model architectures.

\begin{enumerate}
\setlength\itemsep{0pt}
\item \textbf{Coverage screening.} Some variables are missing so some plots have no value at all for that variable, so the whole plot is removed which shrinks the datasets useable rows before any feature set is chosen. This removes 430 plot, from the 4-surveys candidate pool, mostly due to missing data from SoilGrids (413 plots missing), and CEH Topographic Wetness Index (39 plots, mostly
  clustered in one compartment \app{add these missing variable to appendix), plot of them. with examples plot ids.} The 6-survey subset contains no missing values.

\item \textbf{Deduplication.} Deterministic duplicates (variables that are closed-form functions of another candidate) are dropped for example, Near-exact empirical duplicates (Spearman $|\rho| \geq 0.95$) are reduced to one variable, preferring externally measured variables over derived ones because they required fewer additional assumptions.

\item \textbf{Importance ranking.} Three feature-importance signals: univariate Spearman $|\rho|$, permutation importance, and drop-column ablation (the change in held-out $R^2$ when removing a variable), are computed against the target variable and combined by rank-averaging. Correlation ignores context, permutation is biased under correlated predictors, ablation depends on one model/split, so we average across all three balancing out each methods limitations.

\item \textbf{Variance-inflation screening.} To control for multicollinearity, variables are iteratively dropped if their VIF exceeds 5.0. Baseline columns are exempt from removal to anchor all feature sets (Sets 1–5), but still inform VIF calculations for other features.\citep{obrien2007Caution_C5_EVAL}. A
semivariogram, fitted to test-partition residuals only, complements Moran's I with the distance
range over which spatial autocorrelation persists, rather than only whether it is present. %% fix this sentence on the semiveriogram
\end{enumerate}
% (turn above into table to save space if needed)
%% set 1 - 5 table should be after we explain the rsq so it makes more sense why they are different:
\ai{Two separate edits to this block. First: the semivariogram sentence inside step 4 describes an evaluation-time diagnostic, not a feature-screening step -- it doesn't belong here (this is also your own open comment "fix this sentence on the semivariogram", and a second comment at line \S\ref{sec:eval-diagnostics}'s Moran's I paragraph below asks the same placement question). DELETE it from this enumerate and MOVE it, tightened, to \S\ref{sec:eval-diagnostics} beside the Moran's I paragraph, where it belongs alongside the other spatial-autocorrelation diagnostic. Second: acting on your own "(turn above into table to save space if needed)" comment -- yes, do it. DELETE this whole enumerate (4 items) and REPLACE with Table~\ref{tab:feature-screening-steps} below, which keeps every number (430 plots, 413 SoilGrids, 39 CEH TWI, VIF$>5.0$) in a fraction of the space.}

\ai{Replace the preceding enumerate with the following table (real content):}

\begin{table}[!htbp]
\centering
\scriptsize
\renewcommand{\arraystretch}{1.2}
\begin{tabular}{p{2.6cm}p{9.3cm}}
\toprule
Step & What it does \\
\midrule
1. Coverage screening & Plots missing any candidate variable are dropped entirely: 430 plots from
4survey's candidate pool (413 missing SoilGrids, 39 missing CEH Topographic Wetness Index, mostly
clustered in one compartment; plot IDs in Appendix~\ref{app:missing-env}). 6survey has no missing
values. \\
2. Deduplication & Deterministic duplicates (closed-form functions of another candidate) are
dropped; near-exact empirical duplicates (Spearman $|\rho|\geq0.95$) are reduced to one variable,
preferring externally measured over derived variables. \\
3. Importance ranking & Univariate Spearman $|\rho|$, permutation importance, and drop-column
ablation ($\Delta R^2$ when removed) are combined by rank-averaging, balancing each method's
individual weakness (correlation ignores context; permutation is biased under correlated
predictors; ablation depends on one model/split). \\
4. VIF screening & Variables are iteratively dropped if VIF $>5.0$
\citep{obrien2007Caution_C5_EVAL}; baseline columns are exempt from removal but still inform other
variables' VIF. \\
\bottomrule
\end{tabular}
\caption{The four-step screening protocol applied to build each research question's Set1--Set4
(Table~\ref{tab:feature-sets-composition-rq1}).}
\label{tab:feature-screening-steps}
\end{table}

\begin{table}[!htbp]
\centering
\scriptsize
\renewcommand{\arraystretch}{1.2}
\begin{tabular}{p{1cm}p{5.3cm}p{7.2cm}}
\toprule
Set & Selection rule & RQ1 membership (worked example) \\
\midrule
Set1 & Baseline only. & 4 columns: CanopyCover, time\_since\_thinning,
time\_since\_thinning\_missing, recent\_thinning\_5yr. \\
Set2 & Baseline + the top 10 candidates by combined importance rank (Spearman + permutation
importance + drop-column ablation, averaged); a candidate that fails VIF screening is skipped and
replaced by the next-ranked candidate, so the count stays at 10. & 14 columns: Set1 +
gwa\_weibull\_k\_50m, dist\_to\_scpt\_boundary, elevation, tas\_mean, eastness, windward\_topex,
gwa\_weibull\_a\_10m, gwa\_weibull\_a\_50m, slope\_degrees, cpmt\_compactness\_ratio. \\
Set3 & Baseline + the upper half (by combined rank) of the terrain/wind category only -- ranked
within that category, not against every other candidate; VIF-screened. & 15 columns: Set1 +
gwa\_weibull\_k\_50m, elevation, eastness, windward\_topex, gwa\_weibull\_a\_10m,
gwa\_weibull\_a\_50m, slope\_degrees, gwa\_weibull\_k\_10m, solar\_radiation\_index, ceh\_twi,
gwa\_wind\_speed\_10m. \\
Set4 & Set3 + the upper half (by combined rank) of every other category -- climate, soil/site,
spatial-position/edge-effects; VIF-screened. & 18 columns: Set3 + dist\_to\_scpt\_boundary,
tas\_mean, chelsa\_gdd5\_degc, cpmt\_compactness\_ratio. \\
\bottomrule
\end{tabular}
\caption{The four feature sets generated by the screening protocol, using RQ1 as an example. While RQ2 and RQ3 follow the identical protocol, differing candidate pools and target variable for the models lead to different set membership (full lists in Appendix~\ref{app:feature-sets}). Set4 is deliberately the final, richest tier: it is the widest set still curated by evidence-driven relevance ranking, extending Set3's terrain/wind-only ranking to every candidate category.}
\label{tab:feature-sets-composition-rq1}
\end{table}
\ai{CORRECTNESS BUG (not a style choice): the source had two \textbackslash caption commands on this float (one ending "...here is the worked example for RQ1..." and a second, cleaner one immediately after). Only one \textbackslash caption per float is valid -- the first is deleted above, the second (already the better-written one) is kept. Separately, for space: Set2--Set4's "worked example" column currently spells out every one of 10-11 variable names in full, which is exactly the "exhaustive feature configuration" the brief says belongs in the appendix. REPLACE each of those three cells' variable lists with "(e.g. \emph{first two names}; full list Appendix~\ref{app:feature-sets})". This keeps the counts (14/15/18 columns) and one or two concrete examples per set -- enough for a reader to understand what each Set adds -- while cutting roughly two wrapped lines per row.}

\begin{table}[!htbp]
\centering
\scriptsize
\renewcommand{\arraystretch}{1.2}
\begin{tabular}{p{2.1cm}p{0.8cm}p{3.4cm}p{1.5cm}p{1.4cm}p{3.1cm}}
\toprule
Feature set family & Feeds & Model comparison & Comparison role & Attribution role & Purpose \\
\midrule
RQ1 & 18 & DNN and Chapman--Richards-guided network variants (environment-conditioned) & Primary &
No & DNN-vs-guided-network comparison under progressively richer terrain/wind/climate input; no
coefficient-based model reads these lists. \\
RQ2 & 19 & Elastic Net, XGBoost, two-stage mixed-effects model & Secondary & Primary & Which
conditions align with persistent departure from the shared curve. \\
RQ3 & \ct{RQ3's Set4 total column count -- pull from the actual feature manifest, don't guess (not
derivable from the old Set5 total, since Set4 and Set5 were independently VIF-screened, not
nested)} & Elastic Net, XGBoost, GNNWR & Secondary & Primary & Which conditions align with
plot-specific deviation; GNNWR additionally tests whether that association is spatially
varying. \\
\bottomrule
\end{tabular}
\caption{Purpose of the feature sets, each research question uses the identical Set1--Set4 screening protocol, with "Feeds" reflecting the total feature count after deduplication and VIF screening}
\label{tab:feature-set-purpose}
\end{table}
\ai{CORRECTNESS BUG: the source tabular declared 5 columns (p{2.3cm}p{4.2cm}p{1.6cm}p{2cm}p{3.4cm}) but every row supplies 6 fields -- name, Feeds, Model comparison, a Primary/Secondary "comparison role", a Primary/No "attribution role", then Purpose. The header was missing the "Comparison role" label entirely. Fixed above by adding the missing column (both header and column spec) and re-splitting the column widths to fit six columns in scriptsize. No wording changed, this is a rendering fix, not an edit.}

\ai{The pipeline-overview figure that originally sat here has been MOVED to the end of \S1 (Research questions and experimental framework), immediately after Table~\ref{tab:rq-framework} -- see the note there for why (kept and repositioned per your instruction, not cut). Nothing left in this gap; the next section follows directly.}

%------------------------------------------------
\section{RQ1: Top-height prediction models}
\label{sec:rq1-models}

RQ1 evaluates top-height models across four dimensions: baseline predictive accuracy, the use of supplementary environmental data (which can help stakeholders learn more where to invest into further high-quality data acquisition to improve modelling and planning), the benefit of Chapman-Richards physics guidance over unguided networks, and performance generalisation across spatial and temporal splits.

\textbf{Chapman-Richards (CR) curve.} A three-parameter curve
\begin{equation}
H(a) = y_{\max}\left(1-e^{-ka}\right)^{p}
\label{eq:cr}
\end{equation}
predicts top height $H(a)$ at age $a$ using asymptotic height $y_{\max}$, growth rate $k$, and shape parameter $p$. Curves are fitted strictly on each split's training partition via nonlinear least squares (five restarts, lowest sum of squared errors (SSE); Appendix~\ref{app:sse}). Under the pooled $k$-fold split, five separate curves are fitted per cohort (one per fold). The resulting $(y_{\max}, k, p)$ serve as fixed constants in the PINN loss (\S\ref{sec:rq1-pinn-loss}): each PINN model uses strictly its own fold-specific curve so that information from held-out validation or test compartments never leaks into the physics guidance loss.
\ai{Keep this paragraph and Equation~\ref{eq:cr} as-is. ADD one sentence at the end, since it is a genuine, currently-unstated modelling assumption the limitations checklist calls for and it is cheapest to state right where the curve is defined rather than in the cross-cutting limitations section: "The CR form assumes a single monotonic growth trajectory per plot; a plot whose real trajectory is non-monotonic (e.g.\ following disturbance) is not represented by this functional form, and is instead handled by the outlier audit in \S\ref{sec:rq3-target} rather than by the curve itself."}

Three standard supervised models are fit on all RQ1 feature sets (baseline and Sets 1–4) using identical splits. Linear regression fits one.
\textbf{Linear regression.} A single linear relationship was fitted between the predictors and height. Because plots were at least 15 years old in 2008, the observed section of the age--height relationship may appear approximately linear despite the nonlinear full growth trajectory. Linear regression therefore provides a simple baseline, but may perform poorly near the limits of the observed age range. \app{but age distribution, the cut off for early ages (with annotation of unreliable lidar scans) and then an example champan richards curve, for the 4 point example and showing how it aligns with the linear portion of the curve}
\ai{"Linear regression fits one." is a broken leftover sentence fragment (does not parse, duplicates the "\textbf{Linear regression.}" heading right after it) -- DELETE it. Also act on your own \app{} note: it is a good, genuinely space-saving figure request (it replaces the hedge "may appear approximately linear" with an actual visual argument), so convert it to the standard format below instead of leaving it as an open instruction to yourself.}

\ai{FIGURE: Show the training-partition age distribution as a histogram (both cohorts, or 4survey only if space-limited), with the pre-2008 age floor and the associated "unreliable LiDAR scan" exclusion band annotated directly on the x-axis. Overlay a single fitted Chapman-Richards curve (Equation~\ref{eq:cr}) across the same age range, with the observed age window (roughly 15+ years, per the age-floor stated in \S\ref{sec:data-cohorts}) highlighted to show visually that it sits on the curve's near-linear early-to-mid segment. The caption must state that the observed age range covers only part of the full sigmoid trajectory, which is the reason linear regression is a defensible baseline despite the true relationship being nonlinear. This figure replaces the hedge sentence "the observed section of the age-height relationship may appear approximately linear despite the nonlinear full growth trajectory" -- keep only "Linear regression provides a simple baseline but performs poorly near the limits of the observed age range (Figure~\ref{fig:age-linearity})." in the prose.}

\textbf{Random Forest.} Averages bootstrap-aggregated decision trees to capture nonlinear interactions without parametric growth assumptions. This provides a non-parametric inductive bias distinct from the Chapman--Richards and DNN/PINN families, testing whether top-height prediction requires smooth curve constraints or flexible ensemble averaging. (%TODO in background justify this baseline, cite commone sources e.g. lee2018MachineLearningPlotbased observatons doi     = {10.3390/f9050268})
\ai{This paragraph is fine and short, keep as-is. The trailing TODO is a citation placement note for Chapter~\ref{ch:related-work}, not this chapter -- move the note there (or resolve it) rather than leaving a raw TODO comment in a submitted methodology chapter; not something to invent a citation for here.}

\textbf{XGBoost} builds regression trees sequentially, with each tree correcting errors remaining from the preceding ensemble:
\begin{equation}
\hat{y}_i^{(K)} = \sum_{k=1}^{K}\eta f_k(\mathbf{x}_i),
\end{equation}
where $\mathbf{x}_i$ contains the predictors for observation $i$, $f_k$ is the $k$th regression tree, and $\eta$ is the learning rate. Unlike Random Forest, which averages independently fitted trees, this sequential approach can capture complex nonlinear relationships and provides a strong conventional machine-learning comparison for the neural models.

To ensure consistency across all research questions, a single hyperparameter configuration (max depth 4, learning rate 0.04, and up to 500 boosting rounds with early stopping) was applied to every XGBoost model across RQ1–RQ3, rather than tuning each feature set independently. This configuration was carried over from an earlier exploratory fit rather than derived from a systematic grid search. RQ3 used this setup from the outset, whereas RQ1 and RQ2 were initially run on library defaults and subsequently refit under this shared configuration for parity.\sr{rewrd this setnece, too many uses of word conftioufraion} A sensitivity analysis showing the impact of this adjustment on accuracy is provided in Appendix X, with full hyperparameter details in Appendix~\ref{app:hyperparameters}. \app{add to appendix)}
\ai{FACT CHECK, not just a style fix: "RQ1 and RQ2 were initially run on library defaults and subsequently refit under this shared configuration for parity" overstates what has actually happened. As of this pass, RQ2's production attribution model \emph{has} been refit under the shared configuration; RQ1's baseline has NOT been adopted into its production comparison, only run as a separate sensitivity check reported alongside the original (unfixed) fit. Stating both as equally "refit... for parity" would misrepresent RQ1's current status and contradicts the open resolution question you flagged in the results-analysis handover. REPLACE the paragraph (also resolving your own reword note and the vague "Appendix X") with: "To keep hyperparameters comparable across research questions, a single XGBoost configuration (max depth 4, learning rate 0.04, up to 500 boosting rounds with early stopping) is used throughout, carried over from an earlier exploratory fit rather than a systematic grid search. RQ3 has used this configuration from the outset. RQ1's and RQ2's XGBoost models were originally fit on library defaults; RQ2's production models have since been refit under the shared configuration (\S\ref{sec:rq2b}), while RQ1's refit exists only as a sensitivity check reported alongside the original comparison (Appendix~\ref{app:hyperparameters}; outcome in the Results chapter)." This wording is accurate as of today but describes a decision you said is still open (whether to adopt the RQ1 refit into production) -- update this sentence again once that decision is made, don't let it silently go stale.}

\subsection{Neural Network Models}
\label{sec:rq1-pinn-loss}
All DNN and CR-PINN models share an identical core architecture (Figure~\ref{fig:model-architecture}) to ensure that performance differences reflect the loss formulation rather than variations in network capacity. Instead of feeding directly into the base network, environmental covariates enter compact auxiliary sub-networks to adjust the physical curve parameters. This design follows the principles of conditional and variable-coefficient PINNs, where physical constraints or parameters vary with instance-level conditions \citep{kovacs2022Conditional_C4_NN,miao2023VCPINN_C4_NN}. Testing four alternative network sizes confirmed that this shared design performs best across both cohorts, with guided models proving the least sensitive to size changes (Appendix~\ref{app:architecture-screening}).
\ai{RESULTS LEAKAGE: "confirmed that this shared design performs best across both cohorts, with guided models proving the least sensitive to size changes" states the outcome of the architecture-size sensitivity check -- that belongs in the Results chapter, not here (see the results-leakage audit in the accompanying chat response, item 5). REPLACE the last sentence with: "A network-size sensitivity check across four alternative sizes, both cohorts, tests whether this shared size is well-chosen (outcome in the Results chapter; full grids in Appendix~\ref{app:architecture-screening})."}

% FIGURE PLACEHOLDER (model architecture)
% Caption: "Shared DNN/guided-network architecture. The plain data loss (DNN, and the guided
% network's own data term) uses the main trunk's output directly. The guided network adds two
% consistency terms against the Chapman-Richards curve's analytical derivative: an instantaneous
% check via automatic differentiation (physics loss) and a finite-difference check across real
% survey pairs (trajectory loss). In the environment-conditioned variants, terrain/wind features
% feed only the small asymptote (and, in the rate-conditioned variant, rate) sub-networks --
% never the main trunk -- producing a per-plot adjustment to the CR curve's asymptote/rate that
% the physics and trajectory losses are then computed against."
% Content: a block diagram. Left: [Age, no-env features] -> main trunk (3x128, shared) ->
% [predicted height]. Below/beside: [Terrain/wind features] -> small asymptote sub-network (16
% units) -> [y_max adjustment, additive] and, dashed/optional for the rate-conditioned variant,
% -> rate sub-network -> [k adjustment, multiplicative, log-space]. A separate branch shows
% [predicted height] and [age] feeding into an autograd box labelled dH/da, compared against a
% box labelled "CR analytical derivative (frozen y_max, k, p)" to produce the physics loss; a
% second small diagram shows two real survey observations of one plot feeding two forward
% passes, differenced and compared to the same CR derivative at mid-age, to produce the
% trajectory loss. These survey pairs are pre-built once from real consecutive-survey rows before
% training starts (not synthesised per batch) -- draw them as a small fixed table/lookup feeding
% the diagram, not as an on-the-fly sampling step.
\begin{figure}[!htbp]
\centering
\fbox{\parbox{0.9\textwidth}{\centering\vspace{4cm}FIGURE PLACEHOLDER --- Model Architecture Block Diagram\vspace{4cm}}}
\caption{Shared architecture and loss formulation for the standard and Chapman--Richards-guided networks. All variants share a 3-layer predictive trunk ($3 \times 128$). In environment-conditioned models, terrain and wind features feed only auxiliary sub-networks (16 units) to adjust the asymptotic height ($y_{\max}$) and growth rate ($k$). The guided variants augment the standard data loss with: (1) an autograd-derived \emph{physics loss} ($\mathrm{d}H/\mathrm{d}a$) evaluated against the analytical derivative, and (2) a \emph{trajectory loss} computed over consecutive empirical survey pairs.}
\label{fig:model-architecture}
\end{figure}
\ai{KEEP this figure -- it is the clearest case in the chapter of a figure genuinely replacing prose a paragraph could not do as clearly (a dual-loss, dual-branch architecture with a conditioning mechanism). The existing planning comment above is already detailed and correct; it is left as-is (matches the required \ai{FIGURE:...} level of detail already). No further change here beyond what the DNN/CR-PINN prose below is trimmed to reference it rather than re-describe it, noted at those paragraphs.}

The \textbf{Deep neural network} uses three hidden layers of 128 units each, leaky-ReLU activation, age concatenated with every other
input feature. The DNN's loss is plain mean squared error,
\begin{equation}
\mathcal{L}_{\text{data}} = \frac{1}{N}\sum_{i=1}^{N}\left(\hat H_i - H_i\right)^2,
\label{eq:data-loss}
\end{equation}
where $\hat H_i$ is the network's predicted height for row $i$ and $H_i$ is the observed height,
plus an L1 weight penalty ($\lambda_{L1}=10^{-5}$). The environment-conditioned DNN is
architecturally identical to the no-environment DNN; its terrain/wind block is simply concatenated
into the same input, widening that vector rather than adding a second network. This is so that interactions between age and stand variables can be learnt without imposing a separate structure.
\ai{Keep as-is -- this is compact already and Equation~\ref{eq:data-loss} is the actual loss definition, not decoration.}

The \textbf{Chapman-Richards-guided neural network (CR-PINN)} models differ from the baseline DNN only by their loss function. In the baseline no-environment variant, the CR parameters $(y_{\max}, k, p)$ are treated as fixed cohort constants (Equation~\ref{eq:cr}) rather than learned parameters. The model optimises a three-term loss informed by the analytical growth derivative $H'(a) = y_{\max} p k e^{-ka}(1 - e^{-ka})^{p-1}$:

\begin{subequations}
\label{eq:cr-pinn-losses}
\begin{align}
\mathcal{L}_{\text{phys}} &= \frac{1}{N}\sum_{i=1}^{N}\left(\frac{\partial \hat{H}_i}{\partial a_i} - H'(a_i)\right)^2, \label{eq:physics-loss} \\
\mathcal{L}_{\text{traj}} &= \frac{1}{M}\sum_{j=1}^{M}\left(\frac{\hat{H}_{j,\text{later}} - \hat{H}_{j,\text{earlier}}}{\Delta a_j} - H'(\bar{a}_j)\right)^2, \label{eq:trajectory-loss} \\
\mathcal{L} &= \mathcal{L}_{\text{data}} + \lambda_{\text{phys}}\mathcal{L}_{\text{phys}} + \lambda_{\text{traj}}\mathcal{L}_{\text{traj}} + \lambda_{L1}\sum_{\theta}|\theta|. \label{eq:total-loss}
\end{align}
\end{subequations}

Here, the \emph{physics loss} \eqref{eq:physics-loss} checks the model's instantaneous growth rate via automatic differentiation at single points (rescaled by the training-split standard deviation to match real-age units). The \emph{trajectory loss} \eqref{eq:trajectory-loss} checks the predicted height gain across $M$ consecutive real survey visits from the same plot, where $\Delta a_j$ is the time elapsed and $\bar{a}_j$ is the mid-point age (drawn strictly from training rows to avoid leakage). The total objective \eqref{eq:total-loss} combines data loss, both consistency penalties (default $\lambda_{\text{phys}}=\lambda_{\text{traj}}=1.0$; ablation in \S\ref{sec:rq1-ablation}), and $L_1$ regularisation.
\ai{Keep this whole block, including all three sub-equations, unchanged -- this is the chapter's central methodological contribution (the physics-informed loss) and is exactly the kind of equation the brief says must not be cut or exiled to an appendix. ADD one sentence after the paragraph, for the same reason as the CR-curve addition above (a genuinely missing, requested limitation, cheapest stated here): "Because \eqref{eq:physics-loss} and \eqref{eq:trajectory-loss} pull predictions toward the pooled curve's own derivative, this constraint could in principle bias predictions toward a misspecified reference curve rather than the true growth process; the physics-weight ablation (\S\ref{sec:eval-diagnostics}) tests how much the constraint actually shapes predictions rather than assuming it is harmless by construction."}

For the \textbf{Environment-conditioned growth} variants, terrain and wind features $\mathbf{x}^{\text{terrain}}_i$ feed into compact sub-networks (16 units) to produce plot-specific parameter adjustments, leaving the shape parameter $p$ fixed as a global constant:

\begin{subequations}
\label{eq:env-adjustments}
\begin{align}
y_{\max}^{(i)} &= y_{\max} + f_{y}\!\left(\mathbf{x}^{\text{terrain}}_i\right), \label{eq:ymax-adjust} \\
k^{(i)} &= k \cdot \exp\!\left(f_{k}\!\left(\mathbf{x}^{\text{terrain}}_i\right)\right). \label{eq:k-adjust}
\end{align}
\end{subequations}

Here, $f_y$ applies an additive adjustment to the asymptotic ceiling $y_{\max}$ \eqref{eq:ymax-adjust} following site-index modelling principles \citep{socha2023Higher_C1_SS}, while the rate-conditioned variant adds $f_k$ to apply a multiplicative adjustment to the growth rate $k$ \eqref{eq:k-adjust}. Both sub-networks receive only environmental features and are initialised near zero, ensuring training begins from the validated pooled curve before learning local adjustments.
\ai{Keep unchanged -- both equations are single-line, define the actual conditioning mechanism, and Figure~\ref{fig:model-architecture} above already carries the visual walkthrough, so nothing here is redundant with it.}

\paragraph*{Hyperparameter sweeps and ablation.}
\label{sec:rq1-ablation}
To evaluate the true impact of physical guidance, the penalty weights are varied across $\lambda_{\text{phys}} = \lambda_{\text{traj}} \in \{0, 1, 2\}$ and compared against the baseline DNN of identical architecture. The symmetric default ($\lambda = 1.0$) from \eqref{eq:total-loss} was pre-specified before testing instead of being selected after to avoid circular evaluation. Consequently, all guided models in the primary comparison retain this fixed default (full experimental settings and results are detailed in the Results chapter).
\ai{Keep unchanged -- correctly defers the outcome to the Results chapter already, nothing to cut.}

\begin{table}[!htbp]
\scriptsize
\renewcommand{\arraystretch}{1.2}
\begin{tabular}{p{3.2cm}p{3cm}p{2.6cm}p{4.5cm}}
\toprule
Model & Environmental input & Loss terms & Notes \\
\midrule
CR (baseline) & None & N/A (curve fit) & Pooled per cohort, refit per fold for pooled $k$-fold. \\
Linear / random forest / XGBoost & No-environment feature set & Model-specific & XGBoost uses a
fixed configuration, not a validation-selected grid (Appendix~\ref{app:hyperparameters}). \\
DNN (no environment) & None & Data (Eq.~\ref{eq:data-loss}) + L1 & Shared architecture with every
guided variant. \\
DNN (environment-conditioned) & Terrain/wind, concatenated into main input & Data + L1 & Same
network class, wider input only. \\
Guided network (no environment) & None & Data + physics + trajectory + L1 & $y_{\max},k,p$
pooled, frozen. \\
Guided network (terrain-conditioned asymptote) & Terrain/wind, asymptote sub-network only & Data +
physics + trajectory + L1 & Eq.~\ref{eq:ymax-adjust}; $k,p$ pooled. \\
Guided network (terrain-conditioned asymptote and rate) & Terrain/wind, asymptote and rate
sub-networks & Data + physics + trajectory + L1 & Eqs.~\ref{eq:ymax-adjust}--\ref{eq:k-adjust};
$p$ pooled. \\
\bottomrule
\end{tabular}
\caption{Model family comparison for top-height prediction (RQ1). Full hyperparameters:
Appendix~\ref{app:hyperparameters}.}
\label{tab:rq1-models}
\end{table}
\ai{Keep unchanged -- already a compact, correctly-scoped summary table, no results embedded in it.}

%maybe we paraphrase the metrics we will look at here then refer to the later evlaution metrics portion of what these are
\ai{DELETE this comment -- \S\ref{sec:eval-diagnostics} already exists and does exactly this later in the chapter; the comment is an old note-to-self that's now resolved by that section's existence.}

\section{RQ2a: does environment shrink RQ1's own departure from the curve}
RQ2a reuses CR-parameter predictions from RQ1 alongside the fold-matched Chapman--Richards baseline without fitting new models. Deviations from this baseline curve mark plots that diverge from the population average. We test whether environmental covariates reduce these departures, particularly for the most severely over- or under-predicted plots, rather than providing a uniform global improvement.

\paragraph*{Residual reduction}
For each test observation, predictions from the fold-matched CR baseline and the environment-conditioned models (DNN, PINN, and PINN$_k$) are matched by plot and survey year. The residual reduction is computed from the error magnitudes:
\begin{equation}
\text{reduction}_i = \left|r_i^{\text{CR}}\right| - \left|r_i^{\text{model}}\right|,
\label{eq:residual-reduction}
\end{equation}
where a positive value indicates the model predicts closer to the true height than the baseline curve. Evaluating the residual's sign ($r_i = y_i - \hat{y}_i$) alongside \eqref{eq:residual-reduction} identifies whether the model corrects directional over- or under-prediction by the shared curve, rather than merely shrinking error magnitude.
\ai{Keep Equation~\ref{eq:residual-reduction} and this paragraph unchanged -- it is the actual target-construction logic for RQ2a, not restatable in a table without losing precision.}

\paragraph*{Quartile decomposition and coverage.} To separate general model improvements from specific environmental corrections, test observations are partitioned within each fold into equal quartiles by baseline error magnitude $|r^{\text{CR}}|$ (from Q1: smallest error, to Q4: largest). Because this split is purely rank-based, each quartile contains exactly 25\% of the test data; this evaluates performance across standard subgroups rather than focusing on a few extreme outliers. Instead of a single pooled mean, we can separate if environmental conditioning removes a small amount of error across all plots uniformly \textit{or} if conditioning can specifically reduce magnitude of errors on where the shared curve fits worse (Q4), capturing site-specific departures. Mean reductions are evaluated per quartile across all three environment-conditioned models, feature sets (Sets~2--4), and cohorts under spatial cross-validation, with the top-performing configuration verified across multiple random seeds for robustness. %try and define cohort as refring to 4 or 6 in the first place its mentioned in the diss.
\ai{This paragraph's middle sentence ("Instead of a single pooled mean, we can separate if environmental conditioning removes a small amount of error...") restates the same Q1-vs-Q4 hypothesis already given in this section's opening paragraph ("particularly for the most severely over- or under-predicted plots, rather than providing a uniform global improvement") -- DELETE the middle sentence here, it is the same point said twice across two paragraphs. Also resolve your own "define cohort" comment: add "(\S\ref{sec:data-cohorts})" the first time "cohort" appears in this section (in the opening paragraph above) rather than here, then DELETE this trailing comment.}

\section{RQ2b: persistent departure from a shared growth curve}
\label{sec:rq2b}
This asks which environmental/management conditions align with a plot's persistent position above or below the single shared curve, using only the 4survey cohort (6survey's 47 compartments were judged too few for a reliable spatial-block attribution split).
\ai{CORRECTNESS BUG (fixed directly, like the two noted at the top of this file): the source used \textbackslash section* (unnumbered, no ToC entry) here while every other section in the chapter uses \textbackslash section, and had no \textbackslash label at all despite being referenced elsewhere in the chapter as though numbered. An unlabelled starred section breaks \S\ref{} for anyone pointing at it and silently drops it from the table of contents. Changed to a normal numbered \textbackslash section with \textbackslash label\{sec:rq2b\} added; no wording changed.}

\textbf{Target construction.}
For each survey row $i$ in a plot's history, a residual is computed against the pooled baseline curve, and the attribution target $\overline{r}_{\text{plot}}$ is defined as the plot's mean residual across all $T$ observed surveys:

\begin{subequations}
\label{eq:target-construction}
\begin{align}
r_i &= H_i - H(a_i), \label{eq:cr-residual} &
\overline{r}_{\text{plot}} &= \frac{1}{T}\sum_{t=1}^{T} r_{i,t}. \label{eq:mean-cr-residual}
\end{align}
\end{subequations}

Age is excluded from all candidate feature sets because it is the baseline curve's sole predictor in \eqref{eq:cr-residual}, which would introduce circularity into the target in \eqref{eq:mean-cr-residual}.
\ai{Keep unchanged -- this is the target-circularity argument (Age excluded because it built the target), a genuinely load-bearing methodological justification, and the equations are already minimal.}

\paragraph*{Attribution models.}
Three models are fitted on $\overline{r}_{\text{plot}}$ across feature sets (Sets~2--4) under spatial-block cross-validation. Set~1 (baseline only) serves as a targeted ablation to test whether canopy cover signals reflect genuine environmental drivers or prior growth.

\begin{enumerate*}[label=(\roman*)]
\item \textbf{Elastic Net} standardises continuous features, one-hot encodes soil classes, and tunes regularisation parameters ($\ell_1/\ell_2$ ratio and penalty strength) via 5-fold CV on the training partition.
\item \textbf{XGBoost} a separate model from RQ1's, fitted on RQ2's own target
  $\overline{r}_{\text{plot}}$ and RQ2's own candidate features rather than reused from RQ1;
  what carries over is only the hyperparameter recipe, not the fitted model itself (500 boosting
  rounds, maximum tree depth 4, learning rate 0.04, early stopping against a held-out validation fold.) %add to limiations on useing the same hyper parmetrs. for different rq?
\item \textbf{Two-stage mixed-effects model.} Uses a two-stage approximation to a full non-linear mixed-effects model \citep{socha2016Assessment_C2_CR}. Stage~1 isolates persistent plot-level departures ($\overline{r}_{\text{plot}}$) from the frozen baseline curve, and Stage~2 fits a linear mixed-effects model against candidate environmental and management features:
\begin{equation}
\overline{r}_{\text{plot}} = \boldsymbol{\beta}^{\top}\mathbf{x}_{\text{plot}} + u_{\text{cpmt}} + \varepsilon_{\text{plot}},
\label{eq:mixed-effects}
\end{equation}
where $\boldsymbol{\beta}$ represents fixed-effect coefficients on standardised features $\mathbf{x}_{\text{plot}}$, $u_{\text{cpmt}} \sim \mathcal{N}(0, \sigma_u^2)$ is a compartment-level random intercept accounting for spatial clustering, and $\varepsilon_{\text{plot}}$ is the residual error. The model uses random intercepts only; random slopes were omitted because compartment sizes vary widely (from single plots to over a thousand) and no varying-effect hypotheses were posed. Fitted on the training partition just for attribution rather than held-out prediction, coefficients are estimated via restricted maximum likelihood (REML). Nested fixed-effect comparisons use maximum likelihood (ML), as REML likelihoods are not comparable across differing fixed-effects structures \citep{pinheiro2000Mixed_C5_EVAL}. Because $\overline{r}_{\text{plot}}$ is a pre-calculated residual, uncertainty from Stage~1 is not propagated; coefficients are interpreted strictly as statistical attribution of curve departures rather than causal mechanisms.
\end{enumerate*}
%possibly make into a table to save space.
\ai{Considered your own "make into a table" note and decided against it: Elastic Net and XGBoost's descriptions are already 2-3 lines each, and the mixed-effects model needs its equation plus the random-intercepts-only justification in full prose -- a table would either lose that justification or end up no more compact once the equation is accounted for. Keep the enumerate* as-is. DELETE the now-resolved comment. Separately, your own inline note "%add to limitations on using the same hyperparameters for different RQ" is a real, currently-missing limitation -- ADD to the cross-cutting limitations section (\S\ref{sec:limitations}, see the rewritten version in the accompanying chat response) rather than here, then delete this comment.}

Feature importance and SHAP values provide post-hoc interpretability (detailed in \S\ref{sec:eval-diagnostics}). While gain-based importance is extracted directly during tree construction, SHAP values are computed post-fitting across the full dataset using the final model checkpoint. We compare the global ranking of environmental features against the linear mixed-effects coefficients to evaluate whether non-linear tree models and parametric baselines identify consistent environmental drivers, with full evaluation metrics detailed in \S\ref{sec:eval-diagnostics}. -- %change this sentence paragraph to end this section - refer to the actual evaluation metrics use but refer to the next section where these are explained.
\ai{Acting directly on your own instruction here. This paragraph duplicates \S\ref{sec:eval-diagnostics}'s own "Feature attribution (RQ2b, RQ3)" paragraph almost sentence-for-sentence (gain vs.\ SHAP mechanics, RQ2b's cross-model comparison purpose) -- genuine redundancy, not just length. DELETE the whole paragraph and REPLACE with one sentence: "Feature importance and SHAP values provide post-hoc interpretability, compared against the mixed-effects coefficients for cross-model convergence (\S\ref{sec:eval-diagnostics})."}

%=======================================================================================
% From here on (RQ3 onward), the material has not been closely edited yet, so per the brief this
% is rewritten more directly rather than \ai{}-annotated line by line. Larger structural moves are
% marked with plain comments below. Section E of the accompanying chat response has the full
% equation audit; section G has the full results-leakage audit with exact original wording quoted.
%=======================================================================================

\section{RQ3: Plot-specific growth-curve deviation}
\label{sec:rq3-target}
This tests whether a plot's fitted growth ceiling deviates from its static yield-class expectation, and whether local environmental or management conditions explain that deviation. Deviation is estimated in the asymptotic ceiling alone, with the curve's shape parameters held fixed, for three reasons: with few survey observations per plot, fixing shape lets the asymptote be estimated by a stable closed-form linear fit rather than an unstable nonlinear optimisation; the yield-class comparison concerns maximum height potential directly, making the asymptote the parameter of interest; and this follows standard site-index modelling convention, where site quality shifts asymptotic potential while trajectory shape stays constant \citep{socha2021Regional_C2_CR}. By construction the target captures asymptotic height potential alone -- plots with different growth rates or inflection timing that ultimately reach the same ceiling are not treated as divergent (a scope limitation stated once here rather than repeated in \S\ref{sec:limitations}).

\paragraph*{Target construction.}
For each plot, the asymptote $\hat{y}_{\max}^{\text{plot}}$ is estimated from its own repeated height--age surveys via closed-form weighted least squares through the origin, fixing the curve's shape parameters to their Forestry Commission yield-class lookup values. The target is the deviation from the static yield-class asymptote:
\begin{equation}
\Delta y_{\max}^{\text{plot}} = \hat y_{\max}^{\text{plot}} - y_{\max}^{\text{yldc}}.
\label{eq:local-ymax-diff}
\end{equation}
This metric is constructed strictly from height--age records and static tables, with no environmental variables involved.

Plots with extreme $\Delta y_{\max}^{\text{plot}}$ values are audited before modelling against two independent criteria: disturbance flags (clearfelling, measurement anomalies, ambiguous disturbance; exact rules in Appendix~\ref{app:disturbance-rules}) and a trajectory-instability cutoff (top 1\% within-plot height change). This audit's criteria are fixed before any model is fitted; its outcome -- which plots, if any, it flags, and whether they are excluded -- is an empirical result and is reported with RQ3's findings in the Results chapter, not asserted here.
% MOVED TO RESULTS: the original draft stated the audit's outcome directly in this chapter
% ("no checked plots triggered these exclusion criteria, all were retained as genuine empirical
% departures rather than data artifacts"), including a \fc{check this is true} self-flag. That is
% a finding (an empirical claim about what the audit found), not a methodological design choice,
% so it belongs in Results alongside RQ3's other outcomes -- see the results-leakage audit (item 2)
% in the accompanying chat response for the exact original wording and destination subsection.

\paragraph*{Global and spatially varying attribution models.}
Elastic Net and XGBoost are fitted on $\Delta y_{\max}^{\text{plot}}$ using Sets~2--4, following the same protocol as RQ2b (\S\ref{sec:rq2b}) \citep{zou2005Regularization_C4_NN, chen2016XGBoost_C4_NN}. To test whether the environment-deviation relationship varies across the landscape rather than following one global trend, Geographically and Neurally Weighted Regression (GNNWR) is additionally fitted \citep{du2020Geographically_C6_GWR}:
\begin{equation}
\hat{y}_i = \beta_0(s_i) + \sum_j \beta_j(s_i)\,x_{ij},
\label{eq:gnnwr}
\end{equation}
where $s_i$ is plot $i$'s spatial coordinate and each $\beta_j(s_i)$ is a spatially varying coefficient. GNNWR learns how these coefficients vary spatially: distances from each plot to a fixed training reference set feed a neural network that outputs multiplicative weights on a frozen global OLS fit, giving each plot its own local coefficients $\beta_j(s_i)$ (Figure~\ref{fig:rq3-attribution}). This directly tests spatial non-stationarity in the environment-deviation relationship, which a single global model cannot represent by construction.

XGBoost attribution uses gain importance and SHAP values \citep{lundberg2017}. Unlike RQ2b's full-sample SHAP, RQ3's SHAP values are computed out-of-fold across the five spatial-block splits and pooled, so each plot's attribution comes from a model that never trained on its own fold (\S\ref{sec:eval-diagnostics}).

\begin{figure}[!htbp]
\centering
\fbox{\parbox{0.9\textwidth}{\centering\vspace{4cm}FIGURE PLACEHOLDER --- RQ3 target construction
and attribution models\vspace{4cm}}}
\caption{Plot-specific asymptote-deviation target construction (RQ3) and its three attribution
models. Left to right: the closed-form per-plot asymptote fit vs.\ the static yield-class
asymptote (Equation~\ref{eq:local-ymax-diff}) feeds Elastic Net, XGBoost, and GNNWR; GNNWR's box
additionally shows the distance-to-reference-set input, the spatial weighting network, and the
per-predictor weights multiplied against the frozen global OLS coefficients. All three converge on
a spatial-block holdout evaluation; RQ2b's parallel but structurally simpler pipeline (pooled
residual $\to$ Elastic Net/XGBoost/two-stage mixed-effects model) is described in prose and
equations only in \S\ref{sec:rq2b}, not duplicated here, since it has no comparable internal
mechanism to unpack. Annotate: association, not causation.}
\label{fig:rq3-attribution}
\end{figure}
% FIGURE CONTENT (consolidated from the source's separate caption draft / "verified against
% implementation" note list / mermaid sketch / prose fallback -- kept as one structured note
% instead of four overlapping ones):
% - the spatial weighting network's input is a vector of Euclidean distances from the focal plot
%   to every plot in the reference set (MinMax-scaled), NOT raw coordinates or coordinate
%   differences; no k-NN or distance-threshold step exists, "neighbourhood" is learned implicitly
%   from the full distance vector.
% - the network outputs one multiplicative weight per predictor (+ intercept), NOT final
%   coefficients directly; this is multiplied against a separately pre-fitted, frozen global OLS
%   coefficient to give beta_j(s_i). One shared multi-output network produces every coefficient's
%   weight jointly, not one network per coefficient.
% - held-out/test plots reuse the training reference set for their distance vector (not a bespoke
%   neighbour set built from their own location); the reference set is the full, uncapped training
%   partition in the reported runs (confirmed: run with --reference-set-size 0).
% - no geographical kernel or bandwidth parameter exists in the implementation -- do not draw one.
% - the only confirmed outputs are a per-plot coefficient table and a Moran's I residual
%   diagnostic -- label it "per-plot coefficient table", not "coefficient map", unless a
%   map-producing script is separately confirmed.

\section{Model evaluation, uncertainty and interpretation}
\label{sec:eval-diagnostics}
This section states what is measured and why, not what was found (Results chapter).

\paragraph*{Predictive diagnostics (RQ1).}
Every RQ1 model is scored on residuals $e_i = H_i - \hat H_i$ over $N$ rows: MAE $=\frac{1}{N}\sum_i|e_i|$, RMSE $=\sqrt{\frac{1}{N}\sum_i e_i^2}$, $R^2 = 1-\sum_i e_i^2/\sum_i(H_i-\bar H)^2$ (variance explained relative to a mean-only baseline; 1 perfect, 0 matches the mean, negative is worse than the mean on held-out data), and bias $=\frac{1}{N}\sum_i e_i$ (positive is systematic under-prediction). Metrics are further broken down by five age bands to check whether an aggregate score conceals age-specific weaknesses (the maturity gate is applied at plot level, so individual rows below 30 can still appear; a poor $<30$ score reflects a small, non-representative sample, not necessarily worse prediction). Pooled $R^2$ additionally carries a cluster-bootstrap 95\% confidence interval, resampling whole compartments rather than individual rows or plots -- this is RQ1 only; RQ2b and RQ3's own $R^2$ numbers do not currently carry one.
% Converted from a standalone align* block of four metric definitions to inline text: these are
% standard, textbook regression metrics (not a novel or central analytical procedure), so a full
% display equation for each is more space than the definitions need. Kept every definition and
% every interpretive note (age-band caveat, bootstrap-CI scope) exactly as before.

\paragraph*{Spatial autocorrelation (Moran's $I$ and the semivariogram range).}
Global Moran's $I$ for model residuals $e_i = y_i - \hat y_i$ tests whether nearby plots have more similar residuals than distant ones, against a null of spatial randomness:
\begin{equation}
    I = \frac{N}{W} \frac{\sum_{i=1}^{N} \sum_{j=1}^{N} w_{ij} (e_i - \bar{e})(e_j - \bar{e})}{\sum_{i=1}^{N} (e_i - \bar{e})^2}
    \label{eq:morans_i}
\end{equation}
where $\bar e$ is the mean residual, $w_{ij}$ the spatial weight between plots $i$ and $j$, and $W=\sum_{i,j} w_{ij}$ (weighting scheme: Appendix~\ref{app:morans-weights}). Values near zero indicate spatial randomness, positive values residual clustering, negative values dispersion. Moran's $I$ is used two ways: as a post-fit diagnostic on RQ1's raw-height residuals (purely diagnostic, never fed back into feature selection), and as a category-removal check for RQ3 (does removing an environmental category increase residual clustering, currently computed on Set4/4survey only). It has not yet been computed on RQ2b's or RQ3's main attribution-model residuals.

A semivariogram, fitted to test-partition residuals only, complements Moran's $I$ with the distance range over which spatial autocorrelation persists, rather than only whether it is present.
% Relocated from the feature-set-construction section's VIF-screening step, where it did not
% belong methodologically (it is an evaluation-time residual diagnostic, not a step in building
% Set1-4) -- see the \ai{} note at that location above.

\paragraph*{Feature attribution (RQ2b, RQ3).}
Feature importance is evaluated using XGBoost gain and SHAP values \citep{lundberg2017unified}, which decompose individual predictions into signed, additive feature contributions analogous to linear model coefficients. Neither metric is used for feature selection. In RQ2b, SHAP provides a non-linear global attribution view to assess convergence alongside Elastic Net and mixed-effects coefficients; in RQ3, SHAP additionally enables local attribution for individual outlier plots, isolating specific drivers of extreme growth deviations. In both, SHAP values are pooled across the five spatial folds to report mean feature contributions; fold-to-fold SHAP stability is not evaluated.

\paragraph*{Coefficient stability (RQ2b, RQ3).}
Elastic Net coefficients are reported as mean $\pm$ standard deviation across the five spatial folds in both RQ2b and RQ3, to confirm parameter signs and magnitudes remain consistent across spatially distinct partitions rather than reflecting one data split. Attribution models are evaluated across spatial folds without multi-seed sweeps.

\paragraph*{Robustness to stochastic and configuration choices.}
\begin{itemize}
\setlength\itemsep{0pt}
\item \textbf{Training-seed reseed} (RQ1, RQ2a): the RQ1 winner is refit across 4 additional seeds; RQ2a reuses these predictions rather than reseeding separately.
\item \textbf{Architecture-size sensitivity} (RQ1, all three guided-network variants): tests whether the fixed network size affects the headline comparison.
\item \textbf{Physics-weight ablation} (RQ1, PINN and PINN$_k$): tests whether the guided network's advantage, if any, depends on the physics/trajectory loss weighting rather than the mechanism.
\item \textbf{XGBoost hyperparameter sensitivity} (RQ1 baseline): tests whether RQ1's headline comparison depends on the library-default-vs-shared-configuration choice; reported as a sensitivity check alongside the original fit, not yet adopted into RQ1's production comparison (\S\ref{sec:rq1-models}; outcome in the Results chapter). RQ2b's own XGBoost model has already been refit under the shared configuration in production (\S\ref{sec:rq2b}), so it needs no separate sensitivity check here.
\item \textbf{Split-seed reseed} (RQ3, GNNWR, 6survey only): tests whether GNNWR's smaller-cohort result is stable across different fold assignments (\S\ref{sec:limitations}).
\end{itemize}
% Trimmed the last item's original justification clause ("given its wide per-fold variance
% there") -- that pre-states the check's own outcome, which belongs with the other results-leakage
% fixes (see item 4 in the audit).
% Reworded the XGBoost bullet (2026-08-14, per a cross-check from a separate results-analysis
% chat): the original version described RQ1 and RQ2b as equally "being tested," which undersold
% RQ2b's actual status -- RQ2b's refit is the adopted production procedure (matching the corrected
% paragraph in \S\ref{sec:rq1-models}), RQ1's remains sensitivity-check-only. Now says so
% explicitly instead of treating both the same.

\section{Methodological assumptions and limitations}
\label{sec:limitations}
Limitations specific to one model or target are stated beside that method (the CR functional-form assumption in \S\ref{sec:rq1-models}, the physics-loss misspecification risk in \S\ref{sec:rq1-models}, RQ3's asymptote-only target scope in \S\ref{sec:rq3-target}). This section covers the remaining cross-cutting limitations, using the "limited because X, retained because Y" pattern throughout rather than a general disclaimer.

\begin{itemize}
\setlength\itemsep{2pt}
\item \textbf{The spatial buffer is set by convention, not a measured range.} The 60\,m buffer (\S\ref{sec:common-design}) is set from typical plot spacing, not from the residual's own measured spatial-autocorrelation range. A semivariogram-based buffer (\S\ref{sec:eval-diagnostics}) could tighten or loosen this bound, but \textbf{was not used to set the buffer itself because the semivariogram is computed from held-out residuals \emph{after} a split already exists} -- a circular basis for choosing that same split's buffer. Compartment-level holdout remains the primary leakage control; the buffer is a secondary refinement on top of it.
\item \textbf{Feature-set membership is fixed once, not re-derived per fold.} Set1--4 (\S\ref{sec:feature-final}) are ranked and VIF-screened once on the full training partition, not separately per cross-validation fold. A fold-specific ranking would better separate selection variance from model variance, but \textbf{was not used because it would let each of the five folds train on a different feature set}, breaking the direct model-family comparison every research question in this chapter relies on.
\item \textbf{Repeated observations are not modelled as within-plot correlated data.} RQ1's modelling unit is the plot-survey row; the pooled-$R^2$ bootstrap CI (\S\ref{sec:eval-diagnostics}) resamples whole compartments, capturing between-plot spatial clustering, but does not separately model correlation between a single plot's own repeated surveys. \textbf{This is retained because compartment-level resampling is the clustering that matters for RQ1's spatial-generalisation question}; a full hierarchical plot-within-compartment model was judged unnecessary complexity for a predictive-accuracy comparison.
\item \textbf{Environmental variables are point-extracted onto plots smaller than their source resolution.} Terrain/wind/soil rasters (sources and resolutions: \S\ref{sec:data-environmental}) are matched to each plot's centre; several nearby plots can therefore share an identical value from the coarsest sources. \textbf{This reduces the effective variation available to the importance-ranking and VIF steps (\S\ref{sec:feature-final}) for those variables specifically} -- a consequence of the point-extraction choice documented in Chapter~\ref{ch:data}, not a new limitation introduced here.
\item \textbf{Per-plot attribution targets carry unweighted uncertainty.} RQ3's asymptote target (\S\ref{sec:rq3-target}) is fitted from as few as the cohort's minimum number of repeated observations per plot; \textbf{plots with shorter survey histories give a noisier target than plots with longer histories, without this being weighted or flagged downstream}. This compounds for 6survey, whose small compartment count (47) already excluded it from RQ2b's attribution split (\S\ref{sec:rq2b}); GNNWR's own 6survey stability is examined as a separate robustness check (\S\ref{sec:eval-diagnostics}), not assumed here.
\item \textbf{XGBoost's hyperparameter recipe is shared across research questions with different targets.} The same configuration (\S\ref{sec:rq1-models}) is applied to RQ1's height target, RQ2b's mean curve-residual target, and RQ3's asymptote-deviation target, rather than tuned per target. This keeps the three research questions' XGBoost models directly comparable, but \textbf{means none of the three is individually tuned for its own target's scale or noise structure}.
\item \textbf{Association, not causation.} Both attribution designs (RQ2b's fixed-effects regression, RQ3's GNNWR/Elastic Net/XGBoost) report statistical association between environmental/management conditions and curve departure, under the observational design documented in Chapter~\ref{ch:data} (\S\ref{sec:data-strengths-limitations}); \textbf{neither design supports a causal claim}.
\end{itemize}

\section{Reproducibility and implementation}
All split assignments, network weight initialisation, and Elastic Net/XGBoost cross-validation folds use a fixed seed (42) unless explicitly swept as part of a stated robustness check (\S\ref{sec:eval-diagnostics}). Every model reads its feature list from a single generated manifest, so a reported set's exact column membership can always be checked against that source rather than a hand-copied list.

Detailed hyperparameter grids and selected values (Appendix~\ref{app:hyperparameters}), full Set1--Set4 column lists and per-variable screening diagnostics (Appendix~\ref{app:feature-sets}), and secondary sensitivity settings not reported in the main text (dropout-rate/architecture-width sweeps, the narrow-gap temporal split's own results; Appendix~\ref{app:sensitivity}) are kept out of this chapter and pointed to at the relevant section above rather than repeated here.
% Not expanded into a per-choice implementation/limitation/alternative-considered list, despite
% the source's own note requesting this ("add many more implementations... understand the
% limitations and how this affects results, and why we still chose to do it"). That level of
% per-choice tradeoff discussion now lives next to each choice instead (the CR functional-form
% note in \S\ref{sec:rq1-models}, the physics-loss misspecification note in \S\ref{sec:rq1-models},
% the seven cross-cutting items in \S\ref{sec:limitations} above) rather than being duplicated a
% third time here, which the 12-page budget does not have room for.
% DISCLOSED LIMITATION (add to main text/appendix): GNNWR's AIC/AICc are approximated via plain
% feature count rather than the true GWR-corrected effective degrees of freedom, and its F1/F2/F3
% significance tests are unavailable, because the package's exact hat-matrix computation would
% need ~54GB at this project's row counts; R2/RMSE (what GNNWR is actually compared against
% EN/XGBoost on throughout) are computed separately and are unaffected.
